%chktex-file 1 
%chktex-file 8

\documentclass[12pt,a4paper,hidelinks]{article}
\usepackage[british]{babel}

%%MUST STAY HERE
\usepackage[svgnames, table]{xcolor}
%%
\usepackage{algorithm}
\usepackage{algpseudocode}
\usepackage{amsfonts}
\usepackage{amsthm}
\usepackage{amsmath}
\usepackage{array}
\usepackage{amssymb}
\usepackage{arydshln}

\usepackage{booktabs}

\usepackage{eurosym}

\usepackage[T1]{fontenc}
\usepackage{fullpage}
\usepackage{graphicx}

\usepackage[colorlinks = true,
            linkcolor = pdarkerblue,
            urlcolor  = gblue,
            citecolor = gblue,
            anchorcolor = blue]{hyperref}%%Must stay here
\usepackage{caption}
\usepackage{cleveref}
%%
\usepackage{hhline}
\usepackage{listings}

\usepackage{mathrsfs}
\usepackage{mathtools}
\usepackage{minted}
\usepackage{multicol}
\usepackage{multirow}

\usepackage[numbers,sort]{natbib}
\usepackage{navigator}

\usepackage{palatino}
\usepackage{pifont}
\usepackage{pgffor}
\usepackage{pgfplots}
\usepackage{pgfplotstable}
\usepackage{placeins}

\usepackage{sectsty}
\usepackage{siunitx}
\usepackage{subfig}

\usepackage{tabularx}
\usepackage{tikz}
\usepackage{todonotes}

\usepackage{verbatim}

\usepackage{wrapfig}

\usepackage{xinttools}
\usepackage{xcolor}

%%%~~~~~~~~~~~~~~~~~~~~~~~~Fonts~~~~~~~~~~~~~~~~~~~~~~~%%%


%%%~~~~~~~~~~~~~~~~~~~~~~~~Colours~~~~~~~~~~~~~~~~~~~~~%%%
\definecolor{gblue}{HTML}{000099}
\definecolor{pblue}{HTML}{00004d}%{0000FF}
\definecolor{pdarkblue}{HTML}{000080}
\definecolor{pdarkerblue}{HTML}{000066}
%\definecolor{veryblue}{HTML}{00004d}

%~~~~~~~~~~~~~~~~~~~~~~~~~~~~~Tikz~~~~~~~~~~~~~~~~~~~~~~~~~~~~~%
\usetikzlibrary{calc}
\usetikzlibrary{shapes,arrows,arrows.meta}
\usetikzlibrary{automata,positioning}
\usetikzlibrary{spy,backgrounds}
\usetikzlibrary{matrix,positioning,arrows.meta,arrows}
\usetikzlibrary{fadings}
\usetikzlibrary{positioning}
\usetikzlibrary{fit}

\newcommand{\nomefico}{\textbf{nomeFicoDaScegliere}}

\pgfplotsset{compat=1.15}

\captionsetup[figure]{labelfont={color=pdarkblue, bf},textfont={it, small}}

\sectionfont{\bf \Large \color{pblue}}
\subsectionfont{\bf \color{pdarkblue}}
\subsubsectionfont{\bf \color{pdarkerblue}}

%%%%%%%%%%%%%%%%%%%%%%%%%%%%%%%%%%%%%%%%%%%%%%%%%%%%%%%%%%%%%%%%%%%%%%%%%%%%%%%%%%

\begin{document}

\begin{titlepage}
    \centering
    \begin{minipage}{0.22\textwidth}%
        \includegraphics[width=\linewidth]{pics/Cherubino.jpg}
    \end{minipage}\hspace{10pt}
    \begin{minipage}{0.6\textwidth}%
        \flushright
        \large
        \vspace{0.8cm}
        \textsc{\color{pblue}%
        \nomefico\\
        a Bert-based, open domain\\
        question answering system}
        \begin{tikzpicture}%
            \draw[thick, brown] (0.1,0)--(0.99\textwidth,0);%   
        \end{tikzpicture}%
    \end{minipage}%

    
      %{\color{pblue} \textsc{ \Large Human Language Techonlogies}}
    
      %\vspace{0.5cm}
    
      %{\color{pblue} \textsc{ A.Y. 2018-2019}}
    
        \vspace{1.2cm}
    
        Gabriele Barreca, Mario Bonsembiante {\small and} Gemma Martini
        
        \vspace{0.7cm}
        
        University of Pisa
      
        \vspace{1.5cm}
        
        \abstract{Mettere abstract}
    
    %{\large \today}
    \vfill
\end{titlepage}

%%%%%%%%%%%%%%%%%%%%%%%%%%%%%%%%%%%%%%%%%%%%%%%%%%%%%%%%%%%%%%%%%%

\tableofcontents
\thispagestyle{empty}
\newpage
\setcounter{page}{1}

%%%%%%%%%%%%%%%%%%%%%%%%%%%%%%%%%%%%%%%%%%%%%%%%%%%%%%%%%%%%%%%%%%

\section{Introduction}\label{sec:intro}
Question answering (QA) systems can be seen as information retrieval systems which aim is to respond to queries, stated in natural language, by returning short answers or long sentences.
The ``so-called'' \emph{open domain QA task} adds the challenge of understanding if the answer to the selected question may or may not be found in a paragraph, which content has been buried within large text corpora, such as Wikipedia.

\vspace{0.2cm}

Building such systems for practical applications has historically been quite challenging and involved.
In this paper, we present a BERT-based implementation that solves an open domain QA task and we provide some performances study about it.

Moreover, we uploaded a Python implementation on   \href{www.google.com}{this} GitHub repository and we integrated it with a straightforward UI, available at \href{www.gemmamartini.it}{www.gemmamartini.it}.

\nomefico

All references:
\begin{itemize}
  \item kbqa~\cite{kbqa}    
  \item coqa~\cite{coqa}
  \item collobert~\cite{Collobert}
  \item vaswani~\cite{vaswani}
  \item weston1~\cite{weston-tracking}
  \item weston2~\cite{weston-reading}
  \item alberti 2019~\cite{alberti}
  \item kwiatowski 2019~\cite{kwiatowski}
  \item chen~\cite{chen}
  \item liu purple~\cite{liu-purple}
  \item liu yellow~\cite{liu-yellow}
  \item RoBERTa~\cite{roberta}
\end{itemize}

%%%%%%%%%%%%%%%%%%%%%%%%%%%%%%%%%%%%%%%%%%%%%%%%%%%%%%%%%%%%%%%%%%
\section{The architecture}\label{sec:model}
%%%%%%%%%%%%%%%%%%%%%%%%%%%%%%%%%%%%%%%%%%%%%%%%%%%%%%%%%%%%%%%%%%
The broad architecture of \nomefico is shown in \Cref{fig:broad_architecture}.

\begin{figure}[ht!]
    \centering
    \includegraphics[scale=.3]{pics/meanEvents.pdf}
    \caption{Mean of each event happened during one match.}
    \label{fig:meanEvents}
\end{figure}
At this point, we assume that the system receives as input two textual items, the first one is the question and the second one is the paragraph that is supposed to contain the answer

\todo[inline]{short?long?}

In this section the tools used in the project are described.

%%%%%%%%%%%%%%%%%%%%%%%%%%%%%%%%%%%%%%%%%%%%%%%%%%%%%%%%%%%%%%%%%%
\section{Experimental results}\label{sec:experimental_results}
%%%%%%%%%%%%%%%%%%%%%%%%%%%%%%%%%%%%%%%%%%%%%%%%%%%%%%%%%%%%%%%%%%


%%%%%%%%%~~~~~~~~~~~~~~~~~~~~~~~~~~%%%%%%%%%%
\subsection{Hardware}
%%%%%%%%%~~~~~~~~~~~~~~~~~~~~~~~~~~%%%%%%%%%%


%%%%%%%%%~~~~~~~~~~~~~~~~~~~~~~~~~~%%%%%%%%%%
\subsection{Hyper-parameters values}
%%%%%%%%%~~~~~~~~~~~~~~~~~~~~~~~~~~%%%%%%%%%%

%%%%%%%%%%%%%%%%%%%%%%%%%%%%%%%%%%%%%%%%%%%%%%%%%%%%%%%%%%%%%%%%
\section{Future work}\label{sec:future_work}
%%%%%%%%%%%%%%%%%%%%%%%%%%%%%%%%%%%%%%%%%%%%%%%%%%%%%%%%%%%%%%%%
Use a binary classification layer to check if the answer is ``plausible'' (i.e. the meanings in the question are covered in the answer), as done by \cite{Hu2019ReadV} and \cite{Back2020NeurQuRI}.



%%%%%%%%%%%%%%%%%%%%%%%%%%%%%%%%%%%%%%%%%%%%%%%%%%%%%%%%%%%%%%%%
\section{Conclusions}\label{sec:conclusions}
%%%%%%%%%%%%%%%%%%%%%%%%%%%%%%%%%%%%%%%%%%%%%%%%%%%%%%%%%%%%%%%%

%%%%%%%%%%%%%%%%%%%~~~~~~~~~~~~~~~~~~~~~~~~~~%%%%%%%%%%%%%%%%%%%%
\subsection{Boh}

\bibliographystyle{plain}
\bibliography{references}

\end{document} 
